\section{符号说明与基本假设}
\subsection{符号说明}
本文以$X=(g,a,\sigma)$表示可提交调度计划，以$F$表示完成时间，以$Q$表示搬运字节。逻辑张量用于分析跨域通信，片上容量另需区分换出恢复后的物理版本。主要符号见表\ref{tab:symbols}；局部变量在首次使用处定义。
\small
\begin{longtable}{@{}>{\raggedright\arraybackslash}p{0.18\linewidth}@{}>{\raggedright\arraybackslash}p{0.62\linewidth}@{}>{\raggedright\arraybackslash}p{0.10\linewidth}@{}}
\caption{主要符号、含义及单位}\label{tab:symbols}\\
\toprule
符号 & 含义 & 单位\\
\midrule
\endfirsthead
\multicolumn{3}{c}{续表\thetable\quad 主要符号、含义及单位}\\
\toprule
符号 & 含义 & 单位\\
\midrule
\endhead
\midrule
\multicolumn{3}{r}{续下页}\\
\endfoot
\bottomrule
\endlastfoot
$X$ & 可提交的调度方案，$X=(g,a,\sigma)$ & 方案\\
$V$ & 非COPY计算操作集合 & 集合\\
$\mathcal T$ & 逻辑张量集合 & 集合\\
$K$ & 可用核心数 & 核\\
$g(v)$ & 操作$v$所属的子图 & 编号\\
$a(s)$ & 子图$s$所在的核心 & 编号\\
$\sigma_k$ & 核心$k$上的子图执行序列 & 序列\\
$c_v$ & 操作$v$的计算时间 & 周期\\
$b_t$ & 张量$t$的大小 & 字节\\
$p(t)$ & 张量$t$的计算生产者 & 编号\\
$U(t)$ & 张量$t$的计算消费者集合 & 集合\\
$Q_r^0$ & 配置$r$下的切分搬运量 & 字节\\
$Q_r^{\rm spill}$ & 配置$r$下的Spill搬运量 & 字节\\
$D_r$ & 配置$r$下相对整图方案的额外数据搬运量 & 字节\\
$F_r(X)$ & 方案$X$在配置$r$下的Makespan & 周期\\
$L_G$ & 计算依赖图的加权最长路 & 周期\\
$W_M$ & M流水线的总工作量 & 周期\\
$W_V$ & V流水线的总工作量 & 周期\\
$T_{1,i}$ & 第$i$张图的整图单核基准Makespan & 周期\\
$S_{r,i,K}$ & 第$i$张图在配置$r$、$K$核下的加速比 & 1\\
\end{longtable}
\normalsize
其中周期、大小、容量与带宽来自题目配置；方案由求解器选择；事件时刻与Cache状态由评估程序推导。
\subsection{基本假设}
根据题目给出的计算图、硬件参数和评价规则，作如下假设：

（1）附件中的计算图、操作周期、张量大小及依赖关系在一次评价中保持不变；新计划的COPY按题设边界生成。

（2）同一输入图、硬件参数和方案在评估程序中产生确定的Makespan，故可直接比较方案而不引入重复运行噪声。

（3）核内操作调度沿用评估程序的既定规则；本文只决策切图、分核和核内子图顺序。

